\chapter{Architecture du moteur WIDO}
\label{chap:archi}

\section{Vue d'ensemble}

Le moteur \wido{} suit un pipeline en cinq étapes (figure~\ref{fig:pipeline}) :

\begin{enumerate}
    \item \textbf{Parsing} : la requête textuelle est transformée en arbre de décision (AST) par le parseur ;
    \item \textbf{Estimation de cardinalité} : pour chaque clause, un appel API sonde estime le nombre de résultats ;
    \item \textbf{Planification} : le planificateur trie les clauses ET dans l'ordre d'exécution optimal ;
    \item \textbf{Exécution} : le moteur interroge l'API JDM et croise les résultats par jointures sur les IDs ;
    \item \textbf{Présentation} : les résultats sont triés par score, enrichis de preuves, et affichés.
\end{enumerate}

\begin{figure}[H]
    \centering
    % Schéma ASCII simplifié — remplacer par un diagramme TikZ si disponible
    \begin{verbatim}
        Requête utilisateur
              ↓
         [Parseur] → AST
              ↓
      [Cardinalité] → Map{clause → taille}
              ↓
       [Planificateur] → Plan ordonné
              ↓
     [Moteur d'exécution] ↔ [API JDM + Cache]
              ↓
          Résultats
    \end{verbatim}
    \caption{Pipeline de traitement d'une requête WIDO}
    \label{fig:pipeline}
\end{figure}

\section{Technologies utilisées}

\begin{itemize}
    \item \textbf{Environnement d'exécution} : Node.js v14+
    \item \textbf{Serveur HTTP} : Express.js
    \item \textbf{Client HTTP} : Axios (requêtes vers l'API JDM)
    \item \textbf{Cache disque} : fs-extra, fichiers JSON indexés par hash MD5
    \item \textbf{Interface} : HTML5 + CSS3 (glassmorphique) + JavaScript vanilla
\end{itemize}

\section{Structure des fichiers}

\begin{table}[H]
    \centering
    \begin{tabular}{ll}
        \toprule
        Fichier & Rôle \\
        \midrule
        \code{node.js} & Point d'entrée, routage Express, orchestration \\
        \code{parseur.js} & Analyse syntaxique → AST \\
        \code{heuristiques.js} & Planification par complexité et cardinalité \\
        \code{moteurExecution.js} & Exécution, jointures, intersections \\
        \code{jdmApi.js} & Client API JDM + pagination + cache \\
        \code{cacheManager.js} & Cache disque MD5 persistant \\
        \code{benchmark.js} & Tests automatisés + métriques \\
        \code{verify\_logic.js} & Tests de validation du moteur \\
        \bottomrule
    \end{tabular}
    \caption{Fichiers du projet WIDO}
    \label{tab:fichiers}
\end{table}
